\documentclass[11pt,a4paper]{altacv}

\geometry{left=1.5cm,right=1.5cm,top=1.cm,bottom=1.2cm,columnsep=1.5cm}

\usepackage[utf8]{inputenc}
\usepackage[T1]{fontenc}
\usepackage[french]{babel}
\usepackage{paracol}
\usepackage{xcolor}
\usepackage{hyperref}
\usepackage{fontawesome5}
\usepackage{setspace}

\definecolor{SlateGrey}{HTML}{2E2E2E}
\definecolor{LightGrey}{HTML}{666666}
\definecolor{DarkPastelRed}{HTML}{8AC0D9}
\definecolor{PastelRed}{HTML}{00B0DA}
\definecolor{GoldenEarth}{HTML}{B4AD0C}
\colorlet{name}{black}
\colorlet{tagline}{PastelRed}
\colorlet{heading}{DarkPastelRed}
\colorlet{headingrule}{GoldenEarth}
\colorlet{subheading}{PastelRed}
\colorlet{accent}{PastelRed}
\colorlet{emphasis}{SlateGrey}
\colorlet{body}{LightGrey}

\renewcommand{\familydefault}{\sfdefault}

\name{\Large Guillaume BOILEAU}
\tagline{Chef de projet numérique | SI, data, applicatifs et coordination multi-acteurs}
\personalinfo{
  \email{guillaume.boileaupro@gmail.com}
  \phone{+33 6 59 94 85 84}
  \location{Alpes Maritimes, France}
  \linkedin{guillaume-boileau-phd-891b81265}
  \researchgate{Guillaume-Boileau-2}
  \orcid{0000-0002-3576-6968}
}

\setlength{\parskip}{0.75\baselineskip}

\begin{document}
\makecvheader
\vspace{-0.6cm}
\cvsection{Lettre de motivation}
\vspace{-0.4cm}

{\color{accent}\textbf{Objet :}} \textit{Candidature au poste de Chef.fe de projet numérique bâtiment F/H au sein de l'ADEME -- Réf. 2026-4740}

{\color{body}

Madame, Monsieur,

Je vous adresse ma candidature au poste de Chef.fe de projet numérique bâtiment au sein de l'ADEME. Cette opportunité m'intéresse particulièrement car elle correspond à l'évolution que je recherche aujourd'hui : piloter un dispositif numérique à fort enjeu, coordonner des acteurs multiples, structurer des évolutions applicatives et mettre la donnée au service d'une mission d'intérêt général. La perspective de contribuer au développement de l'Observatoire DPE-AUDIT, au croisement des enjeux techniques, réglementaires, data et utilisateurs, constitue pour moi une motivation forte.

Ingénieur de formation et titulaire d'un doctorat, j'ai construit mon parcours dans des environnements complexes où la robustesse des outils, la qualité des livrables, la structuration des besoins et la coordination des contributions sont déterminantes. Au fil de mes expériences, j'ai développé une réelle capacité à piloter des sujets numériques, à organiser les priorités, à suivre des développements, à formaliser des besoins, à accompagner des évolutions techniques et à assurer la cohérence d'ensemble de workflows et d'applicatifs utilisés par plusieurs acteurs.

Chez VEV Platform Services France, j'ai contribué au développement et au suivi de services numériques liés à des flux de données opérationnels dans le domaine de la mobilité électrique. J'y ai travaillé à l'interface entre besoins métier, contraintes techniques, comportements applicatifs et incidents de production. Cette expérience m'a permis de renforcer ma capacité à suivre des services en conditions réelles, à contribuer à l'analyse et à la résolution de dysfonctionnements, à accompagner des évolutions de plateforme, et à intervenir dans une logique de continuité de service, de fiabilité et d'amélioration continue. J'y ai également consolidé une manière de travailler fondée sur la coordination, la réactivité, la priorisation et la recherche de solutions concrètes dans un environnement opérationnel exigeant.

Auparavant, au Laboratoire Lagrange de l'Observatoire de la Côte d'Azur et dans le cadre de travaux menés avec le CNES, j'ai coordonné des développements logiciels et des workflows de données à grande échelle dans un contexte international impliquant plusieurs contributeurs et plusieurs niveaux d'attente. Mon rôle ne se limitait pas au développement technique. Il consistait aussi à structurer les priorités, à clarifier les besoins, à organiser les contributions, à sécuriser la qualité des résultats, à documenter les outils produits et à maintenir une cohérence d'ensemble dans la durée. J'ai notamment conçu une plateforme Python dédiée à l'intégration, à la comparaison et à l'exploitation de données complexes. Ce travail m'a appris à transformer des besoins techniques parfois hétérogènes en un cadre applicatif robuste, lisible et exploitable par différents utilisateurs.

C'est précisément cette articulation entre cadrage, coordination et mise en oeuvre qui me paraît centrale dans le poste proposé par l'ADEME. Piloter une plateforme numérique, organiser et animer une équipe de prestataires, suivre les développements, les recettes et les déploiements, hiérarchiser les évolutions fonctionnelles, prendre en compte les attentes des utilisateurs et contribuer à la qualité de service sont des dimensions dans lesquelles je me reconnais pleinement. J'apprécie les fonctions où il faut à la fois comprendre les enjeux, donner un cadre de travail clair, fluidifier les échanges entre interlocuteurs différents et faire avancer les sujets avec méthode.

Je suis particulièrement sensible à la dimension de traduction d'exigences complexes en solutions concrètes. Sans revendiquer aujourd'hui une spécialisation initiale dans la réglementation du bâtiment, j'ai l'habitude de travailler à partir de contraintes fortes, de besoins évolutifs et d'exigences techniques qu'il faut reformuler, structurer puis rendre opérationnelles dans des outils, des processus ou des workflows. Mon parcours m'a appris à analyser un besoin, à en dégager les implications concrètes, à prioriser les actions nécessaires, puis à accompagner leur mise en oeuvre dans un cadre rigoureux. Cette capacité d'analyse et de transformation me semble directement utile pour un poste où les évolutions de la plateforme doivent rester cohérentes avec des besoins utilisateurs, des objectifs de politique publique et des exigences réglementaires.

Je me reconnais également dans la dimension de pilotage de support et de gestion de demandes complexes. Dans mes expériences précédentes, j'ai souvent dû investiguer des anomalies, suivre des écarts, clarifier des points bloquants, coordonner des actions correctives et maintenir un niveau de lisibilité suffisant pour permettre la prise de décision. J'accorde une grande importance à la qualité du suivi, à la traçabilité des sujets, à la clarté de la documentation et à la capacité de rendre les problèmes compréhensibles et traitables collectivement. Dans un dispositif comme l'Observatoire DPE-AUDIT, cette exigence de méthode et de lisibilité me paraît essentielle.

La contribution à la mise en opendata des données est également un aspect du poste auquel je suis particulièrement sensible. Dans l'ensemble de mon parcours, la qualité, la comparabilité, la structuration et l'exploitation des données ont toujours occupé une place centrale. Je suis convaincu que la valeur d'une plateforme ne repose pas uniquement sur son fonctionnement technique, mais aussi sur sa capacité à produire des données fiables, compréhensibles, bien documentées et réellement utiles aux utilisateurs comme aux décideurs. Dans le cadre d'une mission publique, cette dimension prend encore plus de sens, puisqu'elle participe à la transparence, à l'appropriation des outils et à l'efficacité de l'action menée.

Enfin, je souhaite souligner que mon parcours m'a également permis de développer une réelle capacité de coordination humaine. J'ai supervisé des stagiaires, organisé des contributions d'ingénieurs et de techniciens, structuré des échanges autour d'objectifs communs et maintenu un cadre de travail clair dans des environnements exigeants. J'ai appris à faire avancer des sujets complexes avec des interlocuteurs aux profils variés, à expliciter les priorités, à fluidifier les échanges et à garder une vision d'ensemble sur des projets qui demandent à la fois méthode, autonomie, sens du collectif et constance dans le pilotage.

Rejoindre l'ADEME représenterait pour moi l'opportunité de mettre ces compétences au service d'un projet utile, structurant et directement lié à la transition écologique. Je serais très heureux de pouvoir contribuer au pilotage et au développement de l'Observatoire DPE-AUDIT, en apportant une approche rigoureuse, organisée et orientée vers la qualité du service rendu.

Je serais heureux de pouvoir échanger avec vous afin de vous présenter plus en détail ma motivation et la manière dont je pourrais contribuer efficacement à vos missions.

Je vous prie d'agréer, Madame, Monsieur, l'expression de mes salutations distinguées.

\textbf{Guillaume Boileau}

}

\end{document}